% GENERATED FILE. Edit canonical lessons, then run npm run generate:books.
% canonical-source-hash: fnv1a64:c2ace6fdca749b5c
% canonical-lessons: ZH-C518-qinglang, ZH-C518-yangguang, ZH-W518-yang-sun, ZH-C518-taiyangneng, ZH-C518-dasan, ZH-C518-dawu

\chapter{Sunshine}
\label{ch:zh-sunshine}
\hlchaptermodality{518}

\begin{chapteropening}
\textbf{By the end of this chapter:} \emph{I can talk about clear skies, sunshine, solar power, umbrellas against the sun and fog.}

Complete the last lesson of chapter 518: I can talk about clear skies, sunshine, solar power, umbrellas against the sun and fog.
\end{chapteropening}

\section[qínglǎng]{qínglǎng — clear and cloudless (of the sky)}

\label{lesson:ZH-C518-qinglang}



\begin{quote}
Before the new one: say the Mandarin for starlight, then the Mandarin for twilight.
\end{quote}

\subsection*{You'll want to know: qínglǎng}
\emph{qínglǎng} — \textquotedblleft{}clear and cloudless (of the sky)\textquotedblright{}.

\begin{quote}
\emph{míngtiān tiānqì qínglǎng wǒmen qù gōngyuán} — It'll be clear tomorrow, so we're going to the park.
\end{quote}

Its characters come later; for now it is a word for the ear.

\begin{cousinweb}
\emph{qínglǎng} is built from \emph{qíng} and \emph{lǎng}, both inherited from Middle Chinese, \emph{dzjeng} and \emph{langX} (Guangyun, in Baxter's transcription), which glosses the first as the sky clearing and the second 'bright'.
\end{cousinweb}

\begin{grammarlens}[title={What you've built}]
One more word for saying how it was.
\end{grammarlens}

\subsection*{Your turn}
\begin{itemize}
\raggedright
  \item \emph{Say it:} \emph{qínglǎng}
  \item \emph{Say it:} \emph{qínglǎng}, once more
  \item \emph{Say it:} \emph{míngtiān tiānqì qínglǎng wǒmen qù gōngyuán} — It'll be clear tomorrow, so we're going to the park.
  \item \emph{Recall:} say the Mandarin for starlight, then the Mandarin for twilight, then say \emph{qínglǎng} again
\end{itemize}

\begin{tcolorbox}[breakable,colback=teal!4,colframe=teal!35!black,title={Before you move on}]
What does qínglǎng mean? (\textquotedblleft{}Clear and cloudless (of the sky)\textquotedblright{}.) Say it once more.
\end{tcolorbox}
\section[\texorpdfstring{\zh{阳光} (yángguāng)}{阳光 (yángguāng)}]{\zh{阳光} (yángguāng) — sunshine}

\label{lesson:ZH-C518-yangguang}



\begin{quote}
Before the new one: say the Mandarin for twilight, then the Mandarin for clear and cloudless.
\end{quote}

\subsection*{You'll want to know: \zh{阳光}}
\textbf{\zh{阳光}} — \emph{yángguāng} — \textquotedblleft{}sunshine\textquotedblright{}.

\begin{quote}
\textbf{\zh{这个季节的阳光最好}} \emph{zhège jìjié de yángguāng zuì hǎo} — The sunshine is at its best at this time of year.
\end{quote}

The new character, in pieces: \zh{阳} — six strokes in two pieces: on the left the ear piece, two strokes, as on the right of \zh{那}; on the right \zh{日}, four strokes. The next lesson puts it in your hand.

\begin{cousinweb}
\textbf{\zh{阳光}} (\emph{yángguāng}) is built from \zh{阳} and \zh{光} 'light', both inherited from Middle Chinese, \emph{yang} and \emph{kwang} (Guangyun, in Baxter's transcription). The dictionary glosses the first as the yang of yin and yang and, after an older book, 'high and bright'; the light of the yang is the sun's.
\end{cousinweb}

\begin{grammarlens}[title={What you've built}]
One more word for saying how it was.
\end{grammarlens}

\subsection*{Your turn}
\begin{itemize}
\raggedright
  \item \emph{Say it:} \emph{yángguāng}
  \item \emph{Say it:} \emph{yángguāng}, once more
  \item \emph{Say it:} \emph{zhège jìjié de yángguāng zuì hǎo} — The sunshine is at its best at this time of year.
  \item \emph{Recall:} say the Mandarin for twilight, then the Mandarin for clear and cloudless, then say \emph{yángguāng} again
\end{itemize}

\begin{tcolorbox}[breakable,colback=teal!4,colframe=teal!35!black,title={Before you move on}]
What does \zh{阳光} mean? (\textquotedblleft{}Sunshine\textquotedblright{}.) Say it once more.
\end{tcolorbox}
\section[\texorpdfstring{\zh{阳} (yáng)}{阳 (yáng)}]{\zh{阳} — the ear piece on the left, then \zh{日}}

\label{lesson:ZH-W518-yang-sun}



\begin{quote}
The word first, then the character inside it.

\begin{itemize}
\raggedright
  \item \emph{Recall:} say \emph{yángguāng} — \textbf{R1}, one lesson back, before it decays
\end{itemize}

\begin{quote}
\zh{那} \zh{日}
\end{quote}

\emph{Read it:} each one aloud, with nothing said first — \emph{nà}, \emph{rì}
\end{quote}

\subsection*{Script — the ear piece on the left, then \zh{日}}
\hlblockfigure{\detokenize{figures/ZH-W518-yang-sun-filmstrip.pdf}}{How \zh{阳} is written, stroke by stroke}

\begin{quote}
\zh{阳}
\end{quote}

\textbf{Two pieces, left before right.} Count the pieces, not the strokes.

\textbf{The left piece, the ear, two strokes}, as on the right of \zh{那} and \zh{都}: one stroke that runs right, folds down twice like an ear and hooks, then a long vertical down its left edge.

\textbf{The right piece, \zh{日}, four strokes}, as you wrote it: a vertical, a bar that turns down the right side, a bar across the middle, then a bar closing the foot.

2 + 4 = six.

\subsection*{Your turn}
\begin{itemize}
\raggedright
  \item \emph{Write it:} \zh{阳}, all six, with the model in front of you
  \item \emph{Cover:} write it again from memory, and count the pieces, not the strokes
\end{itemize}

\begin{tcolorbox}[breakable,colback=teal!4,colframe=teal!35!black,title={Before you move on}]
How many pieces does \zh{阳} have? (\textbf{Two: the ear and \zh{日}}.) Which comes first? (\textbf{The ear, on the left}.)

Source: \href{https://raw.githubusercontent.com/chanind/hanzi-writer-data/68d10a4b21150cae5e1ebbd223eed289cf32d90c/data/\%E9\%98\%B3.json}{Hanzi Writer Data \zh{阳}, pinned stroke paths and medians}.
\end{tcolorbox}
\section[\texorpdfstring{\zh{太阳能} (tàiyángnéng)}{太阳能 (tàiyángnéng)}]{\zh{太阳能} (tàiyángnéng) — solar energy; solar power}

\label{lesson:ZH-C518-taiyangneng}



\begin{quote}
Before the new one: say the Mandarin for clear and cloudless, then the Mandarin for sunshine.
\end{quote}

\subsection*{You'll want to know: \zh{太阳能}}
\textbf{\zh{太阳能}} — \emph{tàiyángnéng} — \textquotedblleft{}solar energy; solar power\textquotedblright{}.

\begin{quote}
\textbf{\zh{我们家用了五年太阳能}} \emph{wǒmen jiā yòng le wǔ nián tàiyángnéng} — We've used solar power at home for five years.
\end{quote}

\begin{cousinweb}
\textbf{\zh{太阳能}} (\emph{tàiyángnéng}) is built from \zh{太阳} 'the sun' and \zh{能} 'energy', all inherited from Middle Chinese, \emph{thajH}, \emph{yang} and \emph{nong} (Guangyun, in Baxter's transcription). The dictionary glosses the last 'skilled' and 'able'; the sense 'energy' is modern, and is not given there.
\end{cousinweb}

\begin{grammarlens}[title={What you've built}]
One more word for saying how it was.
\end{grammarlens}

\subsection*{Your turn}
\begin{itemize}
\raggedright
  \item \emph{Say it:} \emph{tàiyángnéng}
  \item \emph{Say it:} \emph{tàiyángnéng}, once more
  \item \emph{Say it:} \emph{wǒmen jiā yòng le wǔ nián tàiyángnéng} — We've used solar power at home for five years.
  \item \emph{Recall:} say the Mandarin for clear and cloudless, then the Mandarin for sunshine, then say \emph{tàiyángnéng} again
\end{itemize}

\begin{tcolorbox}[breakable,colback=teal!4,colframe=teal!35!black,title={Before you move on}]
What does \zh{太阳能} mean? (\textquotedblleft{}Solar energy; solar power\textquotedblright{}.) Say it once more.
\end{tcolorbox}
\section[dǎsǎn]{dǎsǎn — to put up an umbrella}

\label{lesson:ZH-C518-dasan}



\begin{quote}
Before the new one: say the Mandarin for sunshine, then the Mandarin for solar power.
\end{quote}

\subsection*{You'll want to know: dǎsǎn}
\emph{dǎsǎn} — \textquotedblleft{}to put up an umbrella\textquotedblright{}.

\begin{quote}
\emph{wàipó chūmén de shíhou zǒngshì dǎsǎn} — Grandma always puts up an umbrella when she goes out.
\end{quote}

Its characters come later; for now it is a word for the ear.

Two third tones meet in this word: dǎ is said with a rising tone before sǎn, and the spelling keeps both marks.

\begin{cousinweb}
\emph{dǎsǎn} is built from \zh{打} and \emph{sǎn} 'an umbrella', both inherited from Middle Chinese, \emph{taengX} and \emph{sanX} (Guangyun, in Baxter's transcription), which glosses the first 'to strike' and the second 'a canopy'. The vowel of dǎ is not the regular outcome of \emph{taengX}.
\end{cousinweb}

\begin{grammarlens}[title={What you've built}]
One more word for saying how it was.
\end{grammarlens}

\subsection*{Your turn}
\begin{itemize}
\raggedright
  \item \emph{Say it:} \emph{dǎsǎn}
  \item \emph{Say it:} \emph{dǎsǎn}, once more
  \item \emph{Say it:} \emph{wàipó chūmén de shíhou zǒngshì dǎsǎn} — Grandma always puts up an umbrella when she goes out.
  \item \emph{Recall:} say the Mandarin for sunshine, then the Mandarin for solar power, then say \emph{dǎsǎn} again
\end{itemize}

\begin{tcolorbox}[breakable,colback=teal!4,colframe=teal!35!black,title={Before you move on}]
What does dǎsǎn mean? (\textquotedblleft{}To put up an umbrella\textquotedblright{}.) Say it once more.
\end{tcolorbox}
\section[dàwù]{dàwù — thick fog}

\label{lesson:ZH-C518-dawu}



\begin{quote}
Before the new one: say the Mandarin for solar power, then the Mandarin for to put up an umbrella.
\end{quote}

\subsection*{You'll want to know: dàwù}
\emph{dàwù} — \textquotedblleft{}thick fog\textquotedblright{}.

\begin{quote}
\emph{zǎoshang yǒu dàwù wǒmen kàn bu dào shān} — There was thick fog this morning, and we couldn't make out the hills.
\end{quote}

Its characters come later; for now it is a word for the ear.

\begin{cousinweb}
\emph{dàwù} is built from \zh{大} 'big' and \emph{wù} 'fog', both inherited from Middle Chinese, \emph{dajH} and \emph{mjuH} (Guangyun, in Baxter's transcription), which explains the second, after an older book, as what comes when yin and yang are in disorder, and after another, as the earth's breath that the sky does not answer.
\end{cousinweb}

\begin{grammarlens}[title={What you've built}]
One more word for saying how it was.
\end{grammarlens}

\subsection*{Your turn}
\begin{itemize}
\raggedright
  \item \emph{Say it:} \emph{dàwù}
  \item \emph{Say it:} \emph{dàwù}, once more
  \item \emph{Say it:} \emph{zǎoshang yǒu dàwù wǒmen kàn bu dào shān} — There was thick fog this morning, and we couldn't make out the hills.
  \item \emph{Recall:} say the Mandarin for solar power, then the Mandarin for to put up an umbrella, then say \emph{dàwù} again
\end{itemize}

\begin{tcolorbox}[breakable,colback=teal!4,colframe=teal!35!black,title={Before you move on}]
What does dàwù mean? (\textquotedblleft{}Thick fog\textquotedblright{}.) Say it once more.
\end{tcolorbox}
